\documentclass[a4paper,11pt]{article}

\usepackage{v200doc}
\docheader{RadioMaster GX12 Gemini-X}
\hypersetup{pdftitle={Настройка пульта RadioMaster GX12 Dual-Band Gemini-X}}

\begin{document}

\begin{center}
  {\LARGE\bfseries\color{accent} Настройка пульта RadioMaster GX12\\Dual-Band Gemini-X}\\[1.5em]
  \begin{tabular}{rl}
    \textbf{Пульт:}     & RadioMaster GX12 Dual-Band Gemini-X \\
    \textbf{Приёмник:}   & RadioMaster DBR4 Dual-Band Xross Gemini ExpressLRS \\
    \textbf{Прошивка:}   & EdgeTX (пульт) / ExpressLRS (RF-модуль) \\
    \textbf{Режимы:}     & 2.4 GHz, 900 MHz, Gemini-X (Xrossband) \\
  \end{tabular}
\end{center}

\vspace{0.5em}
\noindent\textbf{Источники:}
\begin{itemize}[nosep]
  \item \url{https://www.expresslrs.org/software/gemini}
  \item \url{https://radiomasterrc.com/products/dbr4-dual-band-xross-gemini-expresslrs-receiver}
  \item \url{https://github.com/ExpressLRS/ExpressLRS-Configurator/releases}
  \item \url{https://github.com/EdgeTX/edgetx/releases}
\end{itemize}

\vspace{1em}
\tableofcontents
\newpage

% ============================================================
\section{1. Общее описание}

RadioMaster GX12 --- пульт с двумя встроенными передатчиками ExpressLRS по 1~Вт и складными двухдиапазонными антеннами. Поддерживает диапазоны 2.4~GHz и Sub-G 900~MHz, а также режим \textbf{Gemini-X (Xrossband)} --- одновременная передача на обоих диапазонах для повышения стабильности связи и устойчивости к помехам.

\textbf{Основные характеристики:}
\begin{itemize}[nosep]
  \item \textbf{Каналы:} до 16
  \item \textbf{Экран:} 128$\times$64 монохромный OLED
  \item \textbf{Память:} 512 MB встроенной флеш-памяти
  \item \textbf{Питание:} 7.4~V 2S LiPo или 3.7~V 18650/21700 (до 5000~mAh)
  \item \textbf{Интерфейс:} USB-C (зарядка, прошивка, доступ к хранилищу)
  \item \textbf{Антенны:} встроенные, складные, 2.4~GHz и 900~MHz
\end{itemize}

% ============================================================
\section{2. Обновление прошивки}

\subsection{2.1. Прошивка пульта (EdgeTX) --- EdgeTX Companion}

Прошивку \textbf{самого пульта} (EdgeTX) обновляют через \textbf{EdgeTX Companion}. На Mac~OS~X скачать установщик можно на странице релизов EdgeTX:

\url{https://github.com/EdgeTX/edgetx/releases}

Выбрать сборку \textbf{EdgeTX Companion} для Mac~OS~X (файл вида \texttt{edgetx-companion-<версия>-macos.dmg} или аналог). Подключить пульт по \textbf{USB-C}, запустить Companion, выбрать модель передатчика (GX12) и выполнить прошивку по инструкции в программе.

\subsection{2.2. Прошивка RF-модуля пульта и приёмника (ExpressLRS)}

Прошивку \textbf{RF-модуля} в пульте (ExpressLRS) и приёмника DBR4 обновляют через \textbf{ExpressLRS Configurator}:

\url{https://github.com/ExpressLRS/ExpressLRS-Configurator/releases}

\textbf{Обновление RF-модуля в пульте GX12:}
\begin{enumerate}[nosep]
  \item Установить и запустить ExpressLRS Configurator.
  \item Подключить пульт к компьютеру по \textbf{USB-C}.
  \item Выбрать целевое устройство: \textbf{RadioMaster GX12} (или соответствующий dual-band TX).
  \item Выбрать версию прошивки ExpressLRS и нужные опции (частота, регион, Packet Rate и т.д.).
  \item Нажать \textbf{Build \& Flash} и дождаться завершения прошивки.
\end{enumerate}

\textbf{Обновление приёмника DBR4:}
\begin{enumerate}[nosep]
  \item Подключить DBR4 к компьютеру по USB (через переходник или отладочный порт) либо прошивать по Wi-Fi (режим Wi-Fi на приёмнике).
  \item В Configurator выбрать целевое устройство: \textbf{RadioMaster DBR4} (dual-band RX).
  \item Собрать и прошить прошивку (Build \& Flash или через Wi-Fi).
\end{enumerate}

После обновления пульта и приёмника убедитесь, что версии ExpressLRS совместимы и при необходимости заново выполните связывание (Binding).

% ============================================================
\section{3. Приёмник DBR4 Xross Gemini (двухдиапазонный)}

Приёмник \textbf{RadioMaster DBR4 Dual-Band Xross Gemini} предназначен для работы с пультом GX12 в режиме Gemini-X: принимает сигнал одновременно на 2.4~GHz и 915~MHz, что даёт максимальную стабильность связи.

\subsection{3.1. Характеристики DBR4}

\begin{longtable}{ll}
  \toprule
  \textbf{Параметр} & \textbf{Значение} \\
  \midrule
  RF-чипы          & Semtech LR1121 $\times$ 2 \\
  MCU              & ESP32 \\
  Диапазоны        & 2.404--2.479 GHz / 903.5--926.9 MHz \\
  Скорость пакетов & до 500 Hz (1000 Hz в ExpressLRS v3.4+) \\
  Питание          & 4.5--8.4 V DC \\
  Вес              & 2.70 g без антенн / 5.30 g с антеннами \\
  Размер           & 27.2 $\times$ 27.2 mm \\
  Антенны          & 2$\times$ 2.4 GHz + 2$\times$ 915 MHz (проводные) \\
  Интерфейс        & CRSF (Serial) \\
  Прошивка         & ExpressLRS v3.3.2 (предустановлена) \\
  \bottomrule
\end{longtable}

\subsection{3.2. Связывание (Binding) GX12 и DBR4}

\begin{enumerate}[nosep]
  \item Выключить пульт (GX12).
  \item Подать питание на приёмник, затем \textbf{три раза} быстро включить и выключить питание DBR4.
  \item Дождаться быстрого двойного мигания LED (режим ожидания связывания).
  \item Включить пульт, открыть Lua \textbf{ExpressLRS} $\to$ нажать \textbf{[BIND]}.
  \item После успешного связывания LED на DBR4 горит постоянно.
\end{enumerate}

\subsection{3.3. Режим антенны для пары GX12 + DBR4}

Для полного использования Gemini-X:
\begin{itemize}[nosep]
  \item На \textbf{пульте GX12} в Lua ExpressLRS: \textbf{Antenna Mode} $\to$ \textbf{Gemini}.
  \item На \textbf{приёмнике DBR4} (через Lua при активной телеметрии или ExpressLRS Configurator): \textbf{Antenna Mode} $\to$ \textbf{Gemini}.
\end{itemize}

Оба устройства должны быть в режиме \textbf{Gemini}; тогда пульт работает в Xrossband (2.4 + 900~MHz одновременно).

\subsection{3.4. Подключение к полётному контроллеру (CRSF)}

\begin{itemize}[nosep]
  \item Подключить DBR4 к свободному UART ПК по \textbf{CRSF} (TX приёмника $\to$ RX ПК, RX приёмника $\to$ TX ПК, общая земля).
  \item В ArduPilot: \param{SERIALx\_PROTOCOL} = \textbf{23 (RCIN)} для приёма RC по CRSF.
  \item Скорость: обычно \textbf{400000} или \textbf{115200} (зависит от прошивки ExpressLRS).
\end{itemize}

В комплекте DBR4: приёмник, 2$\times$ антенны 2.4~GHz, 2$\times$ антенны 915~MHz, кабель CRSF, демпферы (4~mm и 3~mm), карта-инструкция.

% ============================================================
\section{4. Режимы передачи ExpressLRS}

\subsection{4.1. Обычный режим (одна антенна)}

Передача только на 2.4~GHz или только на 900~MHz. Используется с приёмниками с одной антенной.

\subsection{4.2. Gemini (однодиапазонный)}

Один диапазон (2.4~GHz или 900~MHz), но пакет одновременно передаётся на двух частотах:
\begin{itemize}[nosep]
  \item 2.4~GHz: разнос частот $\sim$40~MHz
  \item 900~MHz: разнос $\sim$10~MHz
\end{itemize}

Требуется \textbf{diversity-приёмник} (две антенны в одном диапазоне). Даёт более стабильный LQ и лучшую устойчивость к помехам при том же расстоянии.

\subsection{4.3. Gemini-X (Xrossband, двухдиапазонный)}

Передача \textbf{одновременно} на 2.4~GHz и 900~MHz. Один и тот же пакет идёт в обоих диапазонах --- приёмник с поддержкой Gemini-X может принять его с любой стороны. Максимальная устойчивость к помехам и стабильный LQ.

\textbf{Важно:} с приёмником с одной антенной или без Gemini режим Gemini-X на передатчике лучше не использовать: второй <<канал>> будет лишь создавать помехи. В таких случаях выставляйте на пульте режим \textbf{Single Antenna} (Ant1/Ant2) или \textbf{Switch}.

% ============================================================
\section{5. Настройка режима антенны (Lua)}

\begin{enumerate}[nosep]
  \item На пульте открыть меню \textbf{ExpressLRS} (Lua-скрипт).
  \item В разделе настроек передатчика найти \textbf{Antenna Mode} (режим антенны).
  \item Установить режим в соответствии с приёмником:
\end{enumerate}

\begin{longtable}{lll}
  \toprule
  \textbf{Приёмник} & \textbf{Режим на TX (пульте)} & \textbf{Режим на RX} \\
  \midrule
  Одна антенна           & Single (Ant1/Ant2) или Switch & --- \\
  Diversity, без Gemini  & Single / Switch               & --- \\
  Diversity с Gemini     & \textbf{Gemini}               & \textbf{Gemini} \\
  С поддержкой Gemini-X  & \textbf{Gemini} (Gemini-X)    & \textbf{Gemini} \\
  \bottomrule
\end{longtable}

На приёмнике режим антенны тоже задаётся через Lua (при подключении к пульту по телеметрии) или через конфигуратор ExpressLRS.

\textbf{Обязательно:} и на передатчике, и на приёмнике должен быть выбран один и тот же тип режима (оба Gemini или оба Single/Switch).

% ============================================================
\section{6. Model Config Matching (привязка к модели)}

Чтобы для разных моделей автоматически включался нужный режим антенны (например, для одной модели --- Single, для другой --- Gemini):

\begin{enumerate}[nosep]
  \item В EdgeTX: \textbf{Model $\to$ Model Setup} (или аналог в вашей версии).
  \item Включить и настроить \textbf{Model Config Matching} в ExpressLRS.
  \item Назначить, например:
  \begin{itemize}[nosep]
    \item \textbf{Model 5} --- модели с приёмником на одну антенну $\to$ Antenna Mode: Single/Switch.
    \item \textbf{Model 6} --- модели с Gemini-приёмником $\to$ Antenna Mode: Gemini.
  \end{itemize}
\end{enumerate}

Подробнее: \url{https://www.expresslrs.org/model-config-match/}

\end{document}
